\chapter{Proofs for Best-First Search}
\label{annex:p12-heuristic-search}

\begin{annexitem}{annex:p12.hsearch-positive-domain-extension}{thm:p12-hsearch-positive-domain-extension}{Why the comparison proofs have finite witnesses}
\begin{proof}
The positive uniform edge bound and local finiteness give finite
terminating A* and A** runs on solvable inputs, by the proved contour
arguments. Every generated graph at a finite step is finite. Strictly
\(C^*\)-bounded paths have bounded additive cost and bounded depth,
so their entire strict contour is finite too. The strict-contour and
matched-tie proofs therefore apply verbatim: their finite phases, cost
relaxations and returned optimal paths are all still finite.

The consistent-containment proof appends an edge of cost
\(h(v)+\varepsilon>0\) to the original graph. Local finiteness is
preserved, and the new graph has uniform positive bound at least the
minimum of the old bound and that new edge cost. It remains in the
comparison domain. The nonpathological-trade proof instead uses the
finite exposed graph and an explicitly positive appended edge; its
comparison instance is a finite member of the same domain. Thus both
indistinguishability arguments retain their required domain membership.

For normalized path-score forcing, if the formerly used final edge
\(Q-g(P)\) would be zero, replace it by
\(Q+\varepsilon-g(P)>0\), where \(\varepsilon\) is small
enough that \(Q+\varepsilon<C^*\) and
\(Q+\varepsilon<Q_b\). Put the competing direct goal's cost
strictly between \(Q+\varepsilon\) and \(Q_b\).
All path heuristics remain admissible, the old path scores remain
unchanged, and the normalized best-first competitor still selects the
more expensive goal first. This gives the same contradiction using a
finite positive-cost instance. Correctness on the larger domain then
supplies the finite terminating competitor run used in its expansion
and tie arguments.

The weak-compatibility witness still has a fixed solution of cost~11.
Its A* run has a finite bounded-cost contour even with overestimating
advice, and its uniform-cost fallback terminates on every solvable graph
in this domain. The pending-A** wrapper also simulates a finite run;
its early-return certificate only solves shortest paths in a finite
revealed graph. Their strict counterexamples themselves have positive
costs. The costwise proof adds a positive edge and uses the same finite
termination guarantee after the cheaper goal is generated. No arbitrary
infinite graph or nonterminating competitor has been introduced.

Finally suppose \(x>0,y\ge 0\) and \(H(x,y)>x+y\).
Choose \(\varepsilon>0\) and \(q\) with
\(x+y+\varepsilon<q<H(x,y)\). In the two-input counterexample,
give the last edge cost \(y+\varepsilon\), keeping the heuristic
\(h(r)=y\). All three edge costs are positive, the heuristic is
admissible, and the wrong direct goal of cost \(q\) is still
selected. This proves the stated positive-domain bound. For the final
noncontainment assertion, use two root edges of cost~1 and two terminal
edges of cost~1, and estimates~1 at the two non-goal children and zero at
root and goals. The heuristic is consistent. A* policies choosing
opposite tied branches have incomparable expansion sets
\(\{s,a\}\) and \(\{s,b\}\), and both belong to the admissible
normalized path-score class. No successful algorithm can have a subset
of both. This last argument does not assert the corresponding claim
for the stronger costwise class on the positive-cost subdomain.
\end{proof}
\end{annexitem}

\begin{annexitem}{annex:p12.hsearch-positive-type-three}{thm:p12-hsearch-positive-type-three}{A positive-cost first-divergence construction}
\begin{proof}
Compare first expansions when the procedures have the same expanded set
\(D\) and revealed graph \(R\), initially at the root setup.
Orders inside previously matched groups may differ; each completion
preserves each procedure's own earlier observation history. Suppose
the fixed policy is about to expand \(j\) at barrier \(m\).
If \(g(j)<m\), the earlier first-divergence construction already
uses a strictly positive new edge \(m-g(j)\). It gives a proper-subset
challenger whenever the competitor chooses a different next new state;
it also rules out stopping prematurely with an old goal.

It remains to handle \(g(j)=m\), where that edge would be zero.
By the added tie convention, every current frontier state of barrier
\(m\) has \(g=m\), hence heuristic zero. Let \(U_m\) be this
finite nonempty set. All other frontier scores are strictly greater
than \(m\), and equal their endpoint \(g+h\): their witness
ancestors were expanded at barriers at most \(m\). Revealing a
state in \(U_m\) can generate only new path costs strictly greater
than \(m\). Such edges cannot improve the cost of another member
of \(U_m\), or of a previously expanded state whose cost is at
most \(m\).

First observe that the competitor cannot terminate after querying only
a proper subset of \(U_m\) beyond the common revealed set.
Every newly known goal then costs strictly more than \(m\). Append
to an unqueried member a fresh terminal goal by a sufficiently small
positive edge, and terminate all other still unknown lists. The new
goal costs less than the competitor's proposed return. Admissibility is
preserved: old expanded states have \(d_R(s,x)+h(x)\le m\),
newly expanded states of \(U_m\) have heuristic zero, and the
other frontier non-goals can be made dead ends. The triangle inequality
used previously checks every alternative revealed route. The competitor
would miss this cheaper goal. Thus it must either query the whole group
before making an outside query, or make an outside query while a member
remains unqueried. In the first case both procedures expose the same
graph after completing the group; continue the comparison there. In
the second case choose one still unqueried member \(u_0\) immediately
before the first outside query.

Construct a finite instance retaining the old revealed graph and the
original complete successor lists of every member of
\(U_m\setminus\{u_0\}\). Give \(u_0\) only a new goal edge
of cost \(\varepsilon>0\), and give all other unexpanded non-goals
empty successor lists. The retained new edges cannot lower the costs
\(m\) at \(U_m\), so the same triangle argument proves
admissibility of this completion. All newly generated states outside
\(U_m\) have costs, and hence barriers, strictly greater than
\(m\); previously known goals also cost strictly more. There are
finitely many such values. Choose \(\varepsilon\) smaller than
each positive gap to \(m\), ignoring an empty collection of gaps.

The fixed policy selects members of \(U_m\) until it encounters
\(u_0\); no changed list is observable before that point. The new
goal has cost and barrier \(m+\varepsilon\). Before returning it,
the policy expands the remaining states of \(U_m\), all still at
barrier \(m\). Their children are above \(m+\varepsilon\).
Its expansion set on the completion is therefore \(D\cup U_m\).

The competitor's observations remain identical until its first outside
query: all queried lists inside \(U_m\setminus\{u_0\}\) have
been retained.
It therefore pays for such an outside state on the completion too.
Moreover, it cannot finish correctly while leaving any member
\(u\in U_m\) unqueried. To see this, replace that unqueried list
by an edge to a still cheaper fresh goal, of cost
\(0<\varepsilon'<\varepsilon\). This gives a solution of cost
\(m+\varepsilon'\), leaves all its observations unchanged and
retains admissibility by the same inequalities. Any outside non-goal it
has queried is a dead end in the constructed instance and adds no
constraint preventing that completion. Thus admissibility forces the
competitor to query every member of \(U_m\), in addition to the
outside state. Its expansion set properly contains \(D\cup U_m\).
This supplies a proper-subset reverse witness whenever the competitor
leaves a flat group before completing it.

Advance the comparison through identical next queries in the
positive-residual case and through completed flat groups in the
zero-residual case. Group queries neither improve other members nor
reveal a goal at that barrier. A disagreement outside these permitted
reorderings supplies one of the proper-subset witnesses above. A
competitor cannot stop while the fixed policy plans a non-goal query,
by the cheaper-completion arguments. If the fixed policy reaches a
known optimal goal instead, it has no further expansions; the
competitor may only add to the common expansion set. Therefore, unless
a reverse witness occurs, every state expanded by the fixed policy is
expanded by the competitor. A competitor that skips one necessarily
has a reverse witness. This proves type~3.

For no-type~2, perform this comparison with an arbitrary admissible
algorithm on every instance. If any proper-subset witness occurs, the
conclusion follows. Otherwise the fixed policy's expansion set is
contained in the competitor's on every instance. Strict containment
anywhere is itself a type~2 counterexample. If equality holds everywhere,
the positive-cost instance of
Example~\ref{def:p12-hsearch-pathological-subset} supplies the proper
subset \(\{s,b\}\subsetneq\{s,a,b\}\), exactly as in the
earlier no-type~2 proof.

For a solvable locally finite instance with a positive uniform edge
bound, the correctness theorem gives a finite terminating run, so its
first-divergence prefix and each \(U_m\) are still finite. Every
constructed challenger above is a finite positive-cost instance and
belongs to the enlarged domain. The same proof therefore applies there.
\end{proof}
\end{annexitem}

\begin{annexitem}{annex:p12.hsearch-general-closure}{thm:p12-hsearch-general-closure}{The precise fixed-point qualification}
\begin{proof}
Let \(P=(s=v_0,\ldots,v_k)\) be any revealed record. The current
root record has score at most its initial score, because updates are
strict improvements. Suppose inductively that the current record at
\(v_i\) has score at most \(q(P_{s,v_i})\). Continuation-safe
order preservation gives the same inequality after extension by the
next edge of \(P\). That edge has been relaxed after the tail's
last update, and the head retained that proposal or a better one.
Thus the current score at \(v_{i+1}\) is at most
\(q(P_{s,v_{i+1}})\). Induction proves the bound for every revealed
record ending at \(v_k\). The current record itself is one of those
records, so it attains their minimum.

The timing premise cannot be deleted. Use the graph and heuristic in
the proof of Theorem~\ref{thm:p12-hsearch-pending-domination}, but
rank by \(q=g+h\). With ties choosing \(w,x,z,k\) in that order,
the selections begin \(s,w,j,x,z\): \(j\)'s score~5 after
\(w\) is smaller than the remaining scores~10. The proposal~6 from
\(x\) is rejected. After \(z\)'s expansion, \(x\) has
improved to cost~2, but its edge to \(j\) has not yet been relaxed
again. All edges of \(s,z,x,j\) have already been generated, and
its score at \(j\) is~3; the current label at \(j\) is still~5.
The score \(g+h\) is continuation-safe, so this is specifically a
failure of equating generated edges with completed propagation.
\end{proof}
\end{annexitem}

\begin{annexitem}{annex:p12.hsearch-weighted}{thm:p12-hsearch-weighted}{Each weighted and nonlinear bound}
\begin{proof}
On an optimal path, write \(g\) for its prefix cost. Its remaining
cost is \(C^*-g\), so admissibility gives \(h\le C^*-g\).
In the first case, \(p\le 1\) and \(h\ge 0\) give
\[
 q\le g+(1+\varepsilon)h
 \le g+(1+\varepsilon)(C^*-g)
 \le(1+\varepsilon)C^*.
\]
Goal scores equal actual cost. The barrier bound therefore gives the
stated return-quality ratio. Negative values of \(p\) do not
invalidate this upper bound, but can invalidate a separate completeness
condition; no completeness is inferred from the ratio proof.

For the second case, use the stated multiplicative estimate bound
\(h\le\beta(C^*-g)\). Then
\[
 (1-w)g+wh\le(1-w)g+w\beta(C^*-g)
 \le(1-w)C^*
\]
when \(w\beta\le 1-w\). At a goal the score is
\((1-w)C\), strictly increasing in \(C\) because \(w<1\).
The endpoint-maximum certificate proves admissibility. The case \(w=1\)
cannot be included merely when the displayed inequality allows it:
all goal scores would then be zero and strict cost ordering would be lost.

In the third case, monotonicity in the second argument yields
\(H(g,h)\le H(g,C^*-g)\le H(C^*,0)\). The goal normalization
through the strictly increasing function \(C\mapsto H(C,0)\)
again proves admissibility. For the square-root example,
\[
 g^2+(C-g)^2\le C^2
\]
is exactly \(2g(C-g)\ge 0\), proving its required inequality.

Finally put \(r=\phi^{-1}\). Its increase and \(r(0)=0\) are
immediate. To see convexity, concavity gives
\(\phi(\lambda r(a)+(1-\lambda)r(b))\ge\lambda a+(1-\lambda)b\).
Applying the increasing inverse proves the convexity inequality for
\(r\). For \(g,h\ge 0\)
and \(g+h>0\), convexity gives
\[
 r(g)\le\frac{g}{g+h}r(g+h),\qquad
 r(h)\le\frac{h}{g+h}r(g+h).
\]
Adding and applying \(\phi\) proves
\(\phi(r(g)+r(h))\le g+h\). The zero case is immediate,
and \(H(C,0)=C\). The third certificate applies.
\end{proof}
\end{annexitem}

\begin{annexitem}{annex:p12.hsearch-no-type-zero}{thm:p12-hsearch-no-type-zero}{Two incomparable optimal branch choices}
\begin{proof}
Let the root have two distinct successors \(a,b\), each reached by
an edge of cost~1. Each has one zero-cost edge to its own distinct goal.
Set every heuristic to zero. The graph is consistent and has optimum~1.
Uniform-cost search that chooses \(a\) first returns its generated
goal before \(b\), by goal preference, with expansion set
\(\{s,a\}\). The opposite tie policy has set \(\{s,b\}\).
Both are admissible on all finite instances and belong to the normalized
path-score class. Since uniform-cost search remains optimal even under
inadmissible supplied heuristics, they are also costwise compatible.

Any algorithm whose expansion set were contained in both sets would
expand only \(s\). It would not generate either goal and could not
return a revealed solution path. This contradicts success. The example
is pathological, since the non-goal states \(a,b\) have exact
remaining estimates zero. It therefore does not challenge any of the
strict nonpathological-domain theorems.
\end{proof}
\end{annexitem}

\begin{annexitem}{annex:p12.hsearch-general-barrier}{thm:p12-hsearch-general-barrier}{The operational cut and attained minmax}
\begin{proof}
Fix a goal path attaining \(\mu\), truncated at its first goal.
Truncation cannot increase its prefix maximum, so it still attains the
minimum. After the initial expansion, mark each of its non-goal
occurrences processed adequately once a record ending there has been
expanded with score at most that occurrence's source-prefix score.
The first occurrence not yet processed adequately is in OPEN with a
score no greater than its source-prefix score. For path-tree search this
is its own record. For graph merging, the preceding adequate record
generates a no-worse extension by continuation-safe order preservation;
a discarded extension has been replaced by a still better current
record, which either has already been expanded adequately or remains
OPEN. Strict improvement reopens a processed state. The eventual goal
has not been expanded, so the cut cannot disappear before termination.
Its score is at most \(\mu\). Minimum-score selection proves part~1.

Every noninitial prefix in a returned persistent record was itself
selected when it generated the next prefix; the last prefix is the
selected goal. Each such score is at most \(\mu\), so the returned
record has barrier at most \(\mu\). The definition of the minimum
gives the opposite inequality. Its finite maximum therefore equals
\(\mu\), and some noninitial selected prefix attains it. This proves
part~2 without assuming that old ancestor labels automatically equal
shortest values in every previously exposed path.

At the returned goal, \(\psi(C(P_0))=q(P_0)\le\mu\).
If an optimal path has maximum at its goal, its barrier is
\(\psi(C^*)\), so
\[
 \psi(C(P_0))\le\mu\le\psi(C^*).
\]
Strict increase gives \(C(P_0)\le C^*\), and feasibility gives
equality. Applying an inverse where it is defined proves the numerical
bound. Since \(M=\max(q((s)),\mu)\ge\mu\), the full-root bound
is a consequence, not a sharper expansion threshold.

Finally, before the first noninitial selection at \(\mu\), a current
OPEN state with a strictly smaller score cannot remain unselected:
updates only lower its score, and the minimum rule must choose it or a
no-worse replacement before that critical selection. Induction along a
strictly \(\mu\)-bounded path, starting with the initial expansion,
puts each successive state in that situation. A goal cannot intervene
on a path whose entire prefix maximum is strictly below the global
minimum \(\mu\). Thus each required non-goal is expanded before
the first critical selection, proving part~4.
\end{proof}
\end{annexitem}

\begin{annexitem}{annex:p12.hsearch-general-complete}{thm:p12-hsearch-general-complete}{Finiteness of the bounded prefix tree}
\begin{proof}
Consider the rooted tree of all finite records whose noninitial prefix
scores never exceed \(\mu\), stopping each branch at a goal.
This tree is prefix-closed and finitely branching, because the original
graph is locally finite. If it were infinite, its root would have
infinitely many descendants, and one of its finitely many children
would too. Repeatedly choose such a child, using the fixed successor
enumeration to settle choices. This constructs an infinite record chain
whose scores never exceed \(\mu\), contradicting their required
unboundedness above. Therefore the bounded prefix tree is finite.

After the compulsory root, every selected record lies in this tree:
its score is at most \(\mu\), and all its witnessing ancestors
were selected earlier at scores at most \(\mu\). A path-tree
record is selected at most once. In the graph version a state is
reinserted only at a strictly better score, so an identical record
cannot be selected twice. Infinitely many iterations would require
infinitely many distinct selected records in a finite tree. Hence the
algorithm terminates. The maintained cut on an existing solution path
precludes failure. The proof uses only finiteness of this particular
bounded tree; unboundedness on every infinite chain is one sufficient
way to obtain it.
\end{proof}
\end{annexitem}

\begin{annexitem}{annex:p12.hsearch-returned-path}{thm:p12-hsearch-returned-path}{Suffix barriers on the actual returned witness}
\begin{proof}
Each prefix of the returned persistent record was selected before it
generated the next prefix. Fix a prefix ending at \(v_i\). After
that selection and until termination, the usual maintained-cut argument
on the remainder of the returned record gives an OPEN record ending
at some later \(v_j\), with score at most
\(q(P_{s,v_j})\le L_i\). The cut uses an inequality; it does not
require the current label to equal the historical prefix label.

For necessity, take a persistent record at which \(x\) was expanded
and compare it with the returned record. Let \(P_{s,v_i}\) be
their longest common initial segment; if the first record is itself a
prefix of the second, its remaining continuation is empty. Otherwise
all records along its divergent continuation were selected after the
common prefix generated the divergence. During each such selection,
the returned-path suffix cut described above was available. Minimum
selection bounds its score by \(L_i\). Those records provide the
claimed continuation and all of its inequalities.

For sufficiency, choose a later returned-path prefix \(v_j\) at which
the maximum \(L_i\) is attained. The record \(P_{s,v_i}\) has
already been selected when this later record is selected. If the target
state \(x\) has not yet been expanded, the cut on the proposed
continuation has an OPEN record of score no greater than its complete
source-prefix score, hence strictly below \(L_i\). It should be
selected instead of the returned-path record of score \(L_i\), a
contradiction. An empty continuation targets \(v_i\), already
expanded, and requires no separate argument. Continuation-safe merging
preserves the cut by the same one-edge induction as before.
\end{proof}
\end{annexitem}

\begin{annexitem}{annex:p12.hsearch-tree-criterion}{thm:p12-hsearch-tree-criterion}{Both directions of the ordered-tree criterion}
\begin{proof}
First suppose \(k\) is selected. Every noninitial ancestor occurrence
\(n\) on its root path must have been selected as well. If
\(q(n)\ge\ell(n)\), some left side subtree rooted at \(r\),
branching from a strict ancestor of \(n\), contains a goal path
all of whose scores are at most \(\ell(n)\le q(n)\).
Its branching ancestor has already been expanded before \(n\).
Until that goal is selected, there is therefore an OPEN node on this
left goal path with score at most \(q(n)\). Pure-leftmost ties
prefer that node to \(n\). Thus \(n\) cannot be selected before
the left goal, which would terminate the search. This contradiction
proves \(q(n)<\ell(n)\).

If \(q(n)>r(n)\), use a right side subtree attaining \(r(n)\).
Its goal path has every score strictly below \(q(n)\). Its current
OPEN node defeats \(n\) on score alone, giving the same contradiction.
This proves the necessary right inequality. Infinite side minima mean
there is no offending goal path and cause no exceptional finite-score case.

Conversely suppose all inequalities hold. The finite search either
exhausts the tree or terminates at some goal \(t\). Exhaustion selects
every node. Suppose it terminates before selecting \(k\). The goal
cannot be a proper descendant of \(k\), since that would require
\(k\)'s earlier selection, nor a proper ancestor of \(k\), since
such descendants of goals were removed from the search tree. Let
\(a\) be the last common ancestor of \(k,t\), and let \(r\)
be the child of \(a\) on the \(t\)-branch. Put
\(b=\max_{u\in P_{r,t}}q(u)\), and choose a node \(u_0\)
on that path attaining the maximum.

If \(t\) lies to the left, then for every \(n\) on the branch
from just below \(a\) through \(k\),
\[
 q(n)<\ell(n)\le m(r)\le b.
\]
After \(a\)'s expansion, an OPEN node remains on this branch
until \(k\) is selected. Just before \(u_0\)'s selection it
has score strictly less than \(b=q(u_0)\), contradicting minimum
selection. If \(t\) lies to the right, the same argument gives
\(q(n)\le r(n)\le m(r)\le b\); the frontier node is now to
the left of \(u_0\), so equality is also sufficient to defeat
\(u_0\) under pure-leftmost ties. This is again a contradiction.
Both cases prove that \(k\) is selected. The root is compulsory,
and a goal-free tree has all side minima infinite, completing the
remaining boundary cases.
\end{proof}
\end{annexitem}

\begin{annexitem}{annex:p12.hsearch-weak-compatibility}{thm:p12-hsearch-weak-compatibility}{A counterexample to the weaker global premise}
\begin{proof}
Consider the exact revealed pattern
\[
 s\to a:3,\quad s\to b:1,\quad s\to c:1,\quad
 b\to a:1,\quad c\to t:10,
\]
where the displayed lists for \(s,b,c\) are complete, \(t\) is a
goal, and \(a\) has not been expanded. The estimates are
\(h(s)=h(t)=0,h(a)=1,h(b)=20,h(c)=9\). Define an algorithm
\(B\) to inspect this pattern by expanding \(s,b,c\). If any
part of the pattern fails, it uses uniform-cost search; if the exact
pattern holds, it returns the known path through \(c\), costing~11.
The checks are sequential: if a required non-goal state is not generated,
the algorithm falls back before attempting its successor-list query.

On every finite completion of the pattern, the fixed path \(s,c,t\)
keeps an A* frontier score at most~11. Hence A* never expands \(b\),
whose label is~1 and whose score is~21. No alternative route revealed
through \(a\) can improve that label: such a route already costs at
least~3. The only other revealed root branch ends at the terminal goal.
Consequently every new solution A* could discover beyond \(a\)
has a prefix to \(a\) costing~3. Replacing it by \(s,b,a\)
reduces that solution cost by~1. Such a returned solution cannot be
globally optimal. Thus whenever any A* policy is optimal on a completion
of the pattern, its return has the known cost~11, and \(B\)'s return
is optimal too. On all other instances its uniform-cost fallback is
optimal. This proves weak compatibility across the entire domain,
including completions with overestimating advice.

Now take the completion in which \(a\) is a dead end. The heuristic
is admissible, and the optimal path \(s,c,t\) is nonpathological:
its non-goal estimates~0 and~9 are strictly below the true costs~11 and~10.
Every A* policy expands \(a\), whose score~4 is strictly smallest
after the root, but \(B\) skips it. This proves the claimed failure
of strict-contour containment under the weaker premise. It also rules
out repairing that failure merely by choosing different A* ties.
\end{proof}
\end{annexitem}

\begin{annexitem}{annex:p12.hsearch-pending-domination}{thm:p12-hsearch-pending-domination}{A safe early-return wrapper and a strict witness}
\begin{proof}
Fix the following legitimate A** tie policy: after preferring goals,
prefer states never previously expanded to reopened states; within those
groups use a fixed order of state identifiers. Let \(A\) denote that
algorithm. Define \(B\) to simulate \(A\), with the same successive
successor-list queries, but also compute the cheapest path to any known
goal in the finite revealed graph. Immediately before \(A\)'s next
selection, let \(m\) be its minimum OPEN barrier. If the revealed
graph contains a goal path of cost \(K\le m\), return that path;
otherwise make the next simulated query and continue.

On every admissible instance, the frontier proof for A** gives
\(m\le C^*\), while the revealed path gives \(C^*\le K\).
Thus the extra stopping certificate is sound. Finite termination follows
from \(A\)'s termination if the certificate never fires earlier.
Algorithm \(B\) is admissible and its queries form a prefix of
\(A\)'s queries on every instance. Its distinct expansion set is
therefore always a subset of \(A\)'s.

For strictness, use the complete edge list
\[
\begin{split}
 &s\to w:3,\quad s\to x:5,\quad s\to z:1,\quad s\to k:1,\\
 &w\to j:2,\quad x\to j:1,\quad z\to x:1,\quad j\to t:7,
\end{split}
\]
with \(k\) a dead end and \(t\) the sole goal. Put
\[
 h(s)=h(j)=h(t)=0,\quad h(w)=7,\quad h(x)=5,
 \quad h(z)=h(k)=9.
\]
The true remaining costs at \(w,x,z,j\) are~9,8,9,7, so all
estimates are admissible. Choose the fixed order
\(w,x,z,j,k\) among fresh states. All four root successors have
barrier~10. The fresh selections are \(w\), then \(x\), then
\(z\). They give \(g(j)=5\) through \(w\); the proposal~6
through \(x\) was rejected. Expanding \(z\) improves \(x\)
to cost~2 and reopens it at barrier~10, but does not yet relax its cached
edge to \(j\). The fresh-first rule selects \(j\), with its
still pending label~5, and generates the goal at cost~12.

At this point the revealed edges already contain
\(s,z,x,j,t\), costing~10, and the minimum OPEN barrier is~10.
The wrapper \(B\) returns this optimal path. Algorithm \(A\)
instead selects fresh \(k\) before reopened \(x\), and only
afterward propagates the costs~3 at \(j\) and~10 at the goal.
Thus \(B\) skips \(k\), and strict uniform domination follows.
The better route to \(j\) used a rejected non-parent edge from
\(x\); merely updating descendants of the current pointer tree
would not have repaired its pending label.
\end{proof}
\end{annexitem}

\begin{annexitem}{annex:p12.hsearch-saturated-invariant}{thm:p12-hsearch-saturated-invariant}{Cost closure, monotone barriers and a frontier witness}
\begin{proof}
Each accepted label witnesses a simple path, by the strict-improvement
and nonnegative-cycle argument used earlier. The finite revealed graph
therefore admits only finitely many such label improvements. Exhaustive
relaxation of its revealed edges terminates. At a fixed point, following
any finite revealed path gives \(g(v)\) no larger than that path's
cost, whereas its witness proves the reverse inequality against the
minimum. Hence \(g(v)=d_R(s,v)\).

We prove monotonicity of selected barriers by induction. Suppose the next
newly selected non-goal has barrier \(m\), and all earlier selections
have barrier at most \(m\). Every already expanded state's current
\(g+h\) is at most its value at its last selection, and therefore
at most \(m\). A newly improving witness generated by the selected
state has inherited barrier at least \(m\). As its improvement is
propagated through cached expanded states, each new endpoint value
\(g+h\) is still at most \(m\); the inherited barrier stays
exactly \(m\) through that portion of the propagation. At a frontier
state the new barrier is \(\max(m,g+h)\ge m\). Frontier records
not improved already have barrier at least \(m\), by selection of
the minimum. The next selection therefore has barrier at least \(m\),
closing the induction. Initialization has just one record and supplies
the base case. This also proves the asserted bound on \(g+h\) for
previously expanded states, and the selected frontier record itself has
\(g(j)+h(j)\le F(W(j))=m\).

On an admissible instance, a known goal record has barrier equal to its
cost: each heuristic is bounded by that record's own remaining suffix,
and its endpoint score is the total cost. At selection times cost closure
makes this the least revealed goal cost. A non-goal can be selected at
\(m\) only if every known goal costs strictly more, because ties
prefer goals.

For correctness, fix an optimal path and take its first state outside
the expanded set. Its entire prefix has been revealed, so closure gives
an OPEN witness of no greater prefix cost. The suffix argument from
the A** admissibility proof bounds every prefix score of this witness,
and hence its barrier, by \(C^*\). Selection of a goal therefore
returns an optimum. A reachable goal also prevents exhaustion with
failure. On a finite graph only finitely many new states can be expanded.
For a locally finite solvable graph with edge costs at least
\(\delta>0\), selected states have witnesses of cost at most
\(C^*\), hence depth at most \(C^*/\delta\). There are only
finitely many such records. Every finite relaxation phase terminates,
so the whole search terminates in this domain too.
\end{proof}
\end{annexitem}

\begin{annexitem}{annex:p12.hsearch-saturated-type-three}{thm:p12-hsearch-saturated-type-three}{The admissible first-divergence completion}
\begin{proof}
Compare the sequences of first expansions of distinct non-goals by the
fixed \(A^\dagger\) policy and an admissible competitor \(B\).
Suppose they have a common initial sequence with expanded set \(D\)
and identical revealed graph \(R\), after which \(A^\dagger\)
plans to expand \(j\) at barrier \(m\). Construct a finite
instance from \(R\): leave each other unexpanded non-goal with no
outgoing edges, and append a fresh goal \(t\), of heuristic zero,
by an edge from \(j\) of cost
\[
                         c(j,t)=m-d_R(s,j)\ge h(j)\ge 0.
\]
All old revealed edges and heuristic values remain unchanged.

This new instance is admissible. Deleting unknown outgoing edges cannot
shorten a path to an old goal. An expanded state \(x\) that reaches
the new goal must first reach \(j\), and
\[
 d_R(x,j)+m-d_R(s,j)\ge m-d_R(s,x)\ge h(x).
\]
The last inequality is precisely the exposed-graph invariant. The new
edge respects \(j\)'s heuristic as well. Every other frontier
non-goal is terminal and has remaining cost infinity; old goals have
heuristic zero. Thus all cases, including alternative revealed paths to
\(j\), have been checked. The new optimum is \(m\), because
the appended path costs \(m\) and all previously known goal costs
are strictly greater.

The fixed policy makes the same preceding queries and selects \(j\)
on this completion. Its saturated label there is \(d_R(s,j)\),
so the new goal is assigned cost and barrier \(m\). Goal preference
then terminates immediately. Its expansion set is exactly
\(D\cup\{j\}\).

An admissible competitor cannot stop with an old solution after the
common prefix while \(A^\dagger\) plans this new expansion: it
would make the same observations on the completion and miss the cheaper
goal. Thus, if the first-expansion sequences differ, the competitor's
next distinct expansion is some \(k\ne j\), with \(k\notin D\).
It makes that same choice on the completion, since the new edge is still
unrevealed. To remain optimal there it must eventually also expand
\(j\). Its expansion set therefore properly contains
\(D\cup\{j\}\), including the extra state \(k\).

If \(B\) skips a state expanded by \(A^\dagger\) on the
original instance, their finite first-expansion sequences must differ
or one must stop earlier. The premature-stop alternative was just ruled
out. Apply the construction at their first difference. It produces an
admissible finite instance where \(B\) expands a state that the fixed
policy skips, proving the type-3 assertion. The proof requires both
correctness of \(B\) and the permitted zero-cost case of the new
edge; neither is omitted from the statement.
\end{proof}
\end{annexitem}

\begin{annexitem}{annex:p12.hsearch-no-type-two}{thm:p12-hsearch-no-type-two}{Every fixed admissible policy has a pointwise challenger}
\begin{proof}
Fix an admissible algorithm \(B\) and one tie policy of
\(A^\dagger\). If their first-expansion sequences differ on any
admissible finite instance, first handle the case in which
\(A^\dagger\)'s sequence is a proper prefix of \(B\)'s: its
expansion set is already a proper subset on that instance. The opposite
proper-prefix case, where \(B\) stops while \(A^\dagger\)
plans a new non-goal expansion, is precluded by the cheaper-completion
argument. In the remaining case the first-divergence construction in the
preceding proof supplies another instance where
\(E_{A^\dagger}\subsetneq E_B\). Thus \(B\) fails the type-2
test there.

Otherwise their first-expansion sequences, and hence their expansion
sets, agree on every admissible finite instance. Apply
Example~\ref{def:p12-hsearch-pathological-subset}. Cost closure does
not stop \(A^\dagger\) from first expanding \(a\), at barrier~4,
and then \(b\), at barrier~11. Its expansion set is
\(\{s,a,b\}\), while the admissible competitor that inspects
\(b\) first and uses the stated lower-bound certificate expands
only \(\{s,b\}\). The same proper subset therefore challenges
\(B\). Both alternatives produce an admissible challenger, proving
the universal assertion without assuming that an arbitrary \(B\)
is literally a member of the \(A^\dagger\) family.
\end{proof}
\end{annexitem}

\begin{annexitem}{annex:p12.hsearch-bestfirst-forcing}{thm:p12-hsearch-bestfirst-forcing}{A two-goal obstruction and matched boundary choices}
\begin{proof}
Fix such a prefix path \(P:s\leadsto v\), and put
\[
 Q=\max_{u\in P}\{g(P_{s,u})+h(u)\}<C^*.
\]
The initial score is included in this maximum, which is why
\(h(s)<C^*\) is stated. If some noninitial \(b\)-score on this
path were at least \(C^*\), let \(Q_b\) be the maximum of those
noninitial \(b\)-scores. Then \(Q<Q_b\).

Construct a graph consisting of this path, an edge from its end to a new
goal of cost \(Q-g(P)\), and a direct edge from \(s\) to another
new goal of cost \(q=(Q+Q_b)/2\). The first new edge has nonnegative
cost because \(Q\ge g(P)+h(v)\ge g(P)\). Put heuristic zero
at both new goals and preserve all heuristic values on the path. Every
path node satisfies
\(h(u)\le Q-g(P_{s,u})\), so the continuation to the first new
goal respects admissibility; \(h(s)\le Q<q\) also respects the
second goal. Thus the constructed instance is admissible and has optimum
\(Q\).

The path-only restriction preserves every old prefix score \(b\).
After the initial expansion, the direct goal has score \(q\), by
goal normalization. The cheaper solution cannot be reached without
selecting a prefix whose score is \(Q_b>q\). The direct goal is
therefore selected first, yielding a suboptimal return. This contradicts
membership in \(\mathcal A_{\rm bf}\), proving the strict score bound.
The path is simple, so this construction preserves all its state
identities and path data without introducing duplicate-state shortcuts.
Cycle erasure does not weaken the strict-contour application: with
nonnegative edge costs it only decreases the remaining prefix costs and
leaves the state-dependent heuristics unchanged.

The strict prefix bound forces expansions inductively along \(P\).
For a path-tree implementation this follows directly from selecting a
minimum score before any normalized goal of cost \(C^*\). For
continuation-safe graph merging, replacing a prefix by a no-worse
\(b\)-record preserves the upper bound for its next extension. The
strict update/reopening rule maintains a frontier witness until the state
has been expanded. It therefore cannot be skipped. On nonpathological
instances \(E_{A*}=S(I)\), proving universal containment.

For the A** assertion, first take \(h(s)<C^*\). A** expands the
strict contour \(S(I)\) before selecting barrier \(C^*\), as in
the containment proof. The competitor must expand this entire set.
At the boundary, choose A** ties to follow the optimal path returned by
the competitor, processing reopened occurrences when needed. Every
record on that optimal path has barrier at most \(C^*\), so no
off-path non-goal needs selection at the boundary. All its on-path
states were expanded by the competitor. If \(h(s)=C^*\), start
with that tie choice at \(s\); the inherited barrier of every
optimal-path prefix is \(C^*\), so the strict phase is empty and
the same containment follows. The case in which \(s\) is itself
a goal has empty expansion sets.
\end{proof}
\end{annexitem}

\begin{annexitem}{annex:p12.hsearch-two-input-score}{thm:p12-hsearch-two-input-score}{Two alternative solution branches}
\begin{proof}
Suppose \(H(x,y)>x+y\), and choose
\(x+y<q<H(x,y)\). Use distinct states \(s,r,t_1,t_2\), with
edges \(s\to r\) of cost \(x\), \(r\to t_1\) of cost
\(y\), and \(s\to t_2\) of cost \(q\). The last two states
are goals. Put \(h(s)=0,h(r)=y\) and zero heuristics at goals.
The instance is admissible, and its optimum is \(x+y\).
After the start expansion, the non-goal \(r\) has score
\(H(x,y)>q\), whereas the direct goal has normalized score \(q\).
Best-first selection therefore returns the more expensive direct path.
This contradicts admissibility. Zero values of \(x\) or \(y\)
cause no problem in the declared nonnegative finite-graph domain; if
another domain excludes them, the argument only asserts the bound at
permitted inputs.
\end{proof}
\end{annexitem}

\begin{annexitem}{annex:p12.hsearch-tradeoff}{thm:p12-hsearch-tradeoff}{An admissible completion of the exposed graph}
\begin{proof}
Fix an A* run and let \(S=S(I)\), its expansion set on the given
nonpathological instance. Let \(G_e\) contain all generated vertices
and all edges revealed by expansions in \(S\). Give an unexpanded
non-goal vertex no outgoing edges in \(G_e\). Retain each original
goal flag and heuristic. Let \(d_e\) be shortest-path distance in
this finite exposed graph.

At termination, every state of \(S\) has been expanded after its last
strict improvement. Otherwise it would be reopened with score smaller
than its earlier score, hence below \(C^*\), and A* could not yet
select a goal. Each revealed edge out of \(S\) has consequently been
relaxed using its source's final label. These labels are exactly the
shortest costs in \(G_e\): labels witness paths, giving one inequality;
following a shortest exposed path and applying the final relaxation
inequality along every edge gives the other. In particular
\(g(x)=d_e(s,x)\) at every \(x\in S\).

Put
\[
 D=\max_{x\in S}\{d_e(s,x)+h(x)\}<C^*.
\]
The strict inequality includes the initial state by the stated
nonpathological convention. Suppose \(B\) skips \(v\in S\) and
expands no state outside \(S\). Choose \(0<\varepsilon<C^*-D\).
To \(G_e\), append a fresh goal \(t\) with an edge from \(v\)
costing \(D+\varepsilon-d_e(s,v)\). This cost is at least
\(h(v)+\varepsilon>0\).

The new instance is admissible. Deleting unrevealed outgoing edges
cannot decrease any old goal distance. For an expanded state \(x\)
that can reach \(v\), the triangle inequality gives
\[
 d_e(x,v)+D+\varepsilon-d_e(s,v)
 \ge D+\varepsilon-d_e(s,x)>h(x).
\]
Thus the newly added goal also respects its lower bound. An unexpanded
non-goal has no outgoing edge and has remaining cost infinity; an old
goal has normalized heuristic zero. These cover all vertices.

All observations made by \(B\) are unchanged: it expands only old
states of \(S\), never inspects \(v\), and hence sees neither a
deleted outgoing list nor the added edge. It still returns the old
optimum \(C^*\). But the new instance has a goal path costing
\(D+\varepsilon<C^*\), contradicting \(B\)'s admissibility on
every such instance. Therefore skipping a state of \(S\) requires
expanding one outside it, and a proper subset is impossible.

For the final assertion, suppose an admissible algorithm \(U\)
uniformly expands a subset of every admissible competitor's states.
Compare it with A* on any nonpathological admissible instance. Its
expansion set is contained in \(S\), and the trade just proved
prevents it from omitting a member of \(S\). Hence \(E_U=S\).
In Example~\ref{def:p12-hsearch-information-trade}, however, the
admissible competitor skips \(a\in S\), so \(E_U\) is not a
subset of that competitor's expansion set. This contradiction proves
that no such single universally best algorithm exists.
\end{proof}
\end{annexitem}

\begin{annexitem}{annex:p12.hsearch-cost-compatibility}{thm:p12-hsearch-cost-compatibility}{A cheaper return value without an admissibility assumption on the extension}
\begin{proof}
Fix an A* run that expands \(v\) with current label satisfying
\(g(v)+h(v)<C^*\), and suppose \(B\) does not expand \(v\).
Append a fresh goal \(t\), of heuristic zero, with an edge
\(v\to t\) of cost \(h(v)+\varepsilon\), where
\(0<\varepsilon<C^*-g(v)-h(v)\). The modified A* run has the
same observations and selections up to the expansion of \(v\).
That expansion generates a feasible solution of score
\(g(v)+h(v)+\varepsilon<C^*\).

All remaining labels of this goal can only improve. Finite termination
holds even if the extension invalidates heuristic admissibility; while
the goal is in OPEN a larger-score goal cannot be selected. Thus A*
returns a cost strictly less than \(C^*\). Algorithm \(B\), which
never reveals the new successor list, returns the old cost \(C^*\).
This contradicts costwise compatibility. Notice that the proof uses a
comparison of the actual returned costs; it does not infer that A* is
optimal on the modified graph.

For containment up to ties, take the optimal path returned by \(B\)
on the original admissible instance and choose A* ties to advance this
path at score \(C^*\). Its maintained frontier has score at most
\(C^*\); therefore every off-path non-goal expansion has score
strictly below \(C^*\). The argument just proved makes each such
state compulsory for \(B\). Every on-path non-goal was also expanded
by \(B\). This gives the asserted subset. On nonpathological
instances every A* expansion has strict score, except the common initial
state which is already compulsory, so every tie policy has containment.
\end{proof}
\end{annexitem}

\begin{annexitem}{annex:p12.hsearch-pathmax-containment}{thm:p12-hsearch-pathmax-containment}{Matching the search after its strict contour}
\begin{proof}
Both algorithms return cost \(C^*\). Their accepted records use the
same strict \(g\)-improvement test. We first compare their execution
before reaching the score boundary \(C^*\).

Suppose \(h(s)<C^*\). Before A* selects its first record of score
\(C^*\), it expands exactly the states reachable by paths whose
every prefix score is strictly below \(C^*\). Its generated
witnesses cannot cross the boundary before that phase ends. A** has the
same strict paths, because a maximum of finitely many prefix scores is
below \(C^*\) exactly when each prefix score is below it. A better
\(g\)-record whose barrier reaches \(C^*\) cannot prematurely
delete a strict-barrier record: it would have to be generated from an
already selected record of barrier at least \(C^*\), which this
phase has not selected. Conversely a cheaper record generated inside the
strict phase still has a strict barrier whenever the old endpoint score
was strict. Induction along each strict path therefore proves that both
algorithms expand exactly the same finite strict set in this phase.

Every improvement of a strict-set state is processed before the boundary
can be selected. Final relaxation along all exposed edges from that set
makes its labels the least costs achievable through it. This is the
same finite shortest-path closure for both procedures, independently of
their internal order. At every remaining frontier state outside the
strict set, the endpoint score is at least \(C^*\); all proper
prefix scores of its record are smaller. Therefore its A** barrier
equals its A* endpoint score. The two procedures now have identical
frontier costs and identical duplicate labels.

Choose the A** ties to select the same first boundary state as the
specified A* run. If that state is a goal, both terminate. Otherwise,
continue matching A*'s subsequent selections and strict relaxations.
A* selects only endpoint scores at most \(C^*\). Every corresponding
new A** record is generated from a matched record of barrier \(C^*\)
and consequently has barrier \(C^*\) when its endpoint score is
at most \(C^*\). This includes reopened states whose new endpoint
score falls below the boundary: their new record retains the inherited
barrier. Such a record is eligible for the matching A** tie choice.
Records above the boundary need never be selected.

The only obstruction to literal matching is the required preference for
a goal in A** while A* still chooses a non-goal of lower endpoint
score. In that case A** terminates earlier with the already proved
optimal cost and has expanded a subset of the matched run's states.
Otherwise matching continues until both return a goal. If \(h(s)=C^*\),
there is no strict A** phase; start this matching argument at the initial
state. Every A* record subsequently selected has a witnessing prefix
maximum at most \(C^*\), and its inherited initial barrier is
\(C^*\), so the same argument applies. If \(s\) is already a goal,
both expansion sets are empty.

This constructs a matching policy for the specified instance and A*
run. It can be implemented as a tie-choice function on the encountered
histories, with arbitrary legal choices on all other histories; it is
an existence claim, not a rule for selecting optimal ties without
knowing which instance is being appraised. Finally,
Example~\ref{def:p12-hsearch-barrier-tie} gives strict containment:
every A* run expands \(u\), while A**'s required goal preference skips
it. The expansion sets are respectively \(\{s,p,u\}\) and
\(\{s,p\}\).
\end{proof}
\end{annexitem}

\begin{annexitem}{annex:p12.hsearch-consistency-bound}{thm:p12-hsearch-consistency-bound}{Telescoping consistency inequalities}
\begin{proof}
For a finite path \(v=v_0,\ldots,v_k=t\) ending at a goal,
successive edge inequalities give
\[
 h(v_0)\le c(v_0,v_1)+h(v_1)
 \le\cdots\le\sum_{i<k}c(v_i,v_{i+1})+h(t)
 =g(v_0,\ldots,t).
\]
Taking the infimum over goal paths proves \(h(v)\le h^*(v)\).
For a state with no goal path, the right side is \(+\infty\), so
the assertion also holds. On one edge of a start-rooted path, add its
prefix cost to \(h(u)\le c(u,v)+h(v)\), obtaining
\(g(P_{s,u})+h(u)\le g(P_{s,v})+h(v)\).
\end{proof}
\end{annexitem}

\begin{annexitem}{annex:p12.hsearch-frontier}{thm:p12-hsearch-frontier}{The maintained path cut}
\begin{proof}
Write \(P=(v_0=s,v_1,\ldots,v_k)\). Call a non-goal occurrence
\(v_i\) processed sufficiently if that state has been expanded at some
label at most the prefix cost \(g(P_{s,v_i})\). Before termination,
choose the first occurrence on \(P\) not processed sufficiently.
A goal occurrence has not been expanded, so such an occurrence exists.

If it is \(s\), its label zero is in OPEN until its first selection.
Otherwise its preceding occurrence was expanded with an adequate label.
That expansion proposed a label for the chosen state at most its prefix
cost. If the state already had a label at least as good, that label was
retained; otherwise the proposal replaced it and reopened the state if
necessary. Labels never increase. If the chosen state had subsequently
been expanded, it would have been processed sufficiently. Thus it remains
in OPEN with the asserted upper bound. If it is a goal already selected,
the algorithm has terminated with a witness costing no more than the
corresponding prefix, which is the stated exception.

If \(P\) is optimal, its suffix from \(v_i\) is a feasible
continuation. Admissibility gives
\[
 g(v_i)+h(v_i)\le g(P_{s,v_i})+h^*(v_i)
 \le g(P_{s,v_i})+g(P_{v_i,v_k})=C^*.
\]
The argument uses stored label values and persistent witnessing records;
it does not assume that an unprocessed cheaper path has already changed
every descendant's label.
\end{proof}
\end{annexitem}

\begin{annexitem}{annex:p12.hsearch-finite}{thm:p12-hsearch-finite}{Finitely many strict label improvements}
\begin{proof}
Every accepted label is the cost of its stored witnessing path. This is
true at initialization and is preserved because a successful relaxation
extends an existing witness by a real edge of the stated cost.

Every accepted witness is simple. Suppose that an accepted extension to
\(v\) repeated a state already on its preceding witness. At the earlier
occurrence, a label of the corresponding prefix cost had already been
assigned to \(v\). Its current label is no larger. The extension's
cost is at least that old prefix cost, because the intervening cycle has
nonnegative cost. It therefore cannot pass the strict-improvement test.
This contradiction proves simplicity, including the case of a zero-cost
cycle: equality still does not pass the test.

A finite graph has only finitely many simple paths. Hence there are
finitely many successful relaxations and OPEN insertions. Every iteration
removes one current OPEN record; it can be inserted again only after
another successful relaxation. Thus the procedure terminates.

If it returns a goal, its witness is a real path, so it cannot falsely
report success. If it returns failure while a goal path exists, the
frontier-witness lemma supplies a current OPEN state on that path, in
contradiction to the stopping condition. Conversely an unreachable goal
cannot be returned, so finite termination must report failure.
\end{proof}
\end{annexitem}

\begin{annexitem}{annex:p12.hsearch-optimal}{thm:p12-hsearch-optimal}{The lower bound at goal selection}
\begin{proof}
Immediately before any selection, the frontier lemma on an optimal path
supplies an OPEN state with score at most \(C^*\). Selecting a
minimum-score state therefore never selects a score greater than
\(C^*\). If the selected state is a goal \(t\), normalization gives
\(f(t)=g(t)\). Its returned witness is a feasible solution, so
\(C^*\le g(t)\). Together these inequalities yield \(g(t)=C^*\).
This argument does not use consistency and does not suppress reopening.
\end{proof}
\end{annexitem}

\begin{annexitem}{annex:p12.hsearch-infinite}{thm:p12-hsearch-infinite}{A finite bounded-cost contour}
\begin{proof}
Fix one finite solution path \(P\), and put
\[
 M_P=\max_{u\in P}\{g(P_{s,u})+h(u)\}.
\]
This maximum is finite. The same maintained-path-cut argument supplies
an OPEN state of score at most \(M_P\) before termination, whether
or not the heuristic is admissible. Thus each selected record satisfies
\(g\le f\le M_P\), using \(h\ge 0\).

A path of cost at most \(M_P\) has at most
\(\lfloor M_P/\delta\rfloor\) edges. Local finiteness implies,
by induction on this depth bound, that only finitely many such rooted
path records exist. A successful improvement selecting one of them cannot
recur indefinitely, because labels strictly decrease on reinsertion and
there are only finitely many candidate records in this contour. Infinite
execution would therefore require infinitely many selections from a finite
set of possible accepted labels, a contradiction. The frontier lemma also
precludes exhaustion with failure while \(P\) exists.

Finally, any solution path of cost \(C\) bounds the search for the
infimum of solution costs: only finitely many rooted paths cost at most
\(C\), so the minimum among their solution paths exists. If \(h\)
is admissible, the least-cost theorem applies to this attained optimum.
\end{proof}
\end{annexitem}

\begin{annexitem}{annex:p12.hsearch-no-reopen}{thm:p12-hsearch-no-reopen}{No premature settlement with consistency}
\begin{proof}
Suppose a state \(v\) is selected with label \(g(v)>g^*(v)\).
There is a finite path \(P:s\leadsto v\) costing less than that
label, even if the infimum is not initially known to be attained.
The maintained-cut argument gives an OPEN state \(u\) on \(P\)
with label no larger than its prefix cost. To see why the cut has not
disappeared, if every occurrence through \(v\) had already been
expanded at an adequate label, then the nonincreasing current label of
\(v\) could not exceed the cost of \(P\).

Consistency telescopes along the remaining suffix, so
\[
 f(u)\le g(P_{s,u})+g(P_{u,v})+h(v)
       =g(P)+h(v)<g(v)+h(v)=f(v).
\]
This contradicts selection of \(v\). Hence a selected state has
\(g(v)=g^*(v)\); it can never receive a strictly smaller feasible
label later. Reopening is unnecessary under these hypotheses.
\end{proof}
\end{annexitem}

\begin{annexitem}{annex:p12.hsearch-relaxation}{thm:p12-hsearch-relaxation}{The mapped-path lower bound}
\begin{proof}
Replace each edge of any original goal path by its corresponding relaxed
path. Concatenation gives a relaxed goal path of no greater total cost,
because original goals map to relaxed goals. Taking the infimum over
original paths proves \(d_R(\rho(u))\le h^*(u)\).

For an individual edge \((u,v)\), concatenate its relaxed simulation
with an optimal relaxed continuation from \(\rho(v)\), or with a
continuation within arbitrary \(\varepsilon>0\) of the infimum.
Then
\[
 d_R(\rho(u))\le c(u,v)+d_R(\rho(v))+\varepsilon.
\]
Letting \(\varepsilon\) tend to zero gives consistency. At a goal,
the empty relaxed continuation costs zero; nonnegative costs force the
normalized value zero. A suboptimal relaxed solution supplies an upper
bound on \(d_R\), which need not be a lower bound on the original
optimum, so it cannot be substituted in this proof without another
inequality.
\end{proof}
\end{annexitem}

\begin{annexitem}{annex:p12.hsearch-sure-set}{thm:p12-hsearch-sure-set}{Strictly bounded paths and the equality boundary}
\begin{proof}
Let \(P:s\leadsto v\) have all its noninitial prefix scores below
\(C^*\). The initial state is expanded unless it is already a goal.
Inductively, expansion of an adequate label at one prefix generates an
adequate label at the next state. That state's score is below \(C^*\)
and remains so under every subsequent improvement. If it has not already
been expanded, it must be selected before a goal of score \(C^*\).
Finite termination and optimality ensure that such a goal is eventually
selected. Thus every state on \(P\), including \(v\), is expanded.

Now fix a nonpathological optimal path. Until its goal is generated,
the maintained cut on that path contains a non-goal of score strictly
below \(C^*\). After its goal is generated, a minimum-score selection
of score \(C^*\) returns a goal because goals have priority. Therefore
every expanded noninitial record has score strictly below \(C^*\).
Its witnessing record was generated from an expanded parent record;
recursing through these persistent witnesses shows that every earlier
noninitial prefix also had score below \(C^*\). The state consequently
belongs to \(S(I)\). Both containments prove equality of the sets,
independently of ties among equally scored non-goals.
\end{proof}
\end{annexitem}

\begin{annexitem}{annex:p12.hsearch-consistent-efficiency}{thm:p12-hsearch-consistent-efficiency}{An indistinguishable cheaper completion}
\begin{proof}
The initial state is compulsory for every nontrivial solution, so consider
a noninitial \(v\in S(I)\). Some path to \(v\) has prefix score
at \(v\) strictly less than \(C^*\); hence
\(g^*(v)+h(v)<C^*\). Choose
\[
 0<\varepsilon<C^*-g^*(v)-h(v).
\]
Append a fresh goal \(t\) by a single edge from \(v\) costing
\(h(v)+\varepsilon\), and put \(h(t)=0\). Leave all old
vertices, edges, goal flags and estimates unchanged. Consistency holds
on every old edge and on the new one. Thus the extended instance is
admissible, and it has a solution of cost
\(g^*(v)+h(v)+\varepsilon<C^*\).

If an equal-information admissible algorithm \(B\) does not expand
\(v\) in the original instance, none of its observations distinguishes
the extended instance: the only new edge lies in an uninspected successor
list. It therefore returns the same old cost \(C^*\), contradicting
its guarantee on all admissible instances. This proves that \(B\)
must expand \(v\). The argument uses the specified information oracle;
it would not apply if the entire graph or an exact complete vertex/goal
list were already supplied to \(B\).

On a nonpathological consistent instance, A* expands exactly \(S(I)\),
so universal containment follows. For a general consistent instance, fix
the cheapest solution path returned by \(B\), erasing zero-cost loops
if needed. Every proper prefix state of that path was expanded by \(B\).
Choose A* ties to advance that path whenever the minimum score equals
\(C^*\). Consistency gives least-cost labels at selections and
nondecreasing prefix scores. Every off-path selection of score below
\(C^*\) belongs to \(S(I)\) and is compulsory for \(B\).
No off-path non-goal with score \(C^*\) needs to be selected: the
next unprocessed path state is available with score at most \(C^*\).
If another goal is selected earlier at the same optimum, no further
non-goal expansions are added. Thus this tie policy expands a subset
of \(E_B(I)\). The choice depends on the competitor's returned path,
as the theorem explicitly permits.
\end{proof}
\end{annexitem}

\begin{annexitem}{annex:p12.hsearch-pathmax}{thm:p12-hsearch-pathmax}{A path barrier with cost-based duplicate handling}
\begin{proof}
The witness and strict-improvement invariants from the finite termination
proof depend on the relaxation test, not on the ordering of OPEN.
They therefore hold for A** as well. In particular, finite graphs yield
only finitely many accepted labels and hence termination; a reachable
goal prevents failure by the same path-cut argument.

Fix an optimal path \(P\). Its maintained cut supplies an OPEN state
\(v\) with witnessing record \(W\) satisfying
\(g(W)\le g(P_{s,v})\). For every ancestor occurrence \(u\) in
\(W\), the remainder of \(W\) from \(u\) to \(v\), followed
by the suffix of \(P\), is a feasible goal continuation. Consequently
\[
 g(W_{s,u})+h(u)
 \le g(W)+g(P_{v,t})\le C^*.
\]
Taking the maximum over all occurrences in \(W\) proves
\(F(W)\le C^*\), without assuming that \(W\) literally coincides
with the optimal path's prefix.

For any record ending at a goal, admissibility applied to its own suffixes
gives every prefix score at most its total cost, while the last score
equals that cost. Thus \(F(W(t))=g(W(t))\). A minimum-barrier
selection cannot return a goal of cost greater than \(C^*\), since
the frontier witness has barrier at most \(C^*\). Feasibility gives
the reverse inequality.

On a locally finite graph with costs at least \(\delta>0\), selected
records satisfy \(g\le F\le C^*\), because \(h\ge 0\). There
are only finitely many rooted paths of at most
\(\lfloor C^*/\delta\rfloor\) edges. The finite-contour argument
therefore establishes termination in this case also. Notice that this
proof uses cost-based duplicate updates; replacing that test by a
comparison of the barriers is a different algorithm and has not been
justified here.
\end{proof}
\end{annexitem}
